% figures/w02-pseudo-derivative.svg — 미분에는 이득의 천장이 없다.
%
%   bash _tools/tikz2svg.sh figures/src/w02-pseudo-derivative.tex figures/w02-pseudo-derivative.svg
%
% 크기응답을 TeX 이 직접 계산한다. 좌표를 미리 구워 넣지 않으므로
% Kd·N 을 바꾸면 곡선이 따라 움직인다.
%
%   ideal   |Kd jw|            = Kd w                 (끝없이 상승)
%   pseudo  |Kd N jw/(jw+N)|   = Kd N w / sqrt(w^2+N^2)
%                              -> Kd w   (w << N)
%                              -> Kd N   (w >> N)      평탄
\documentclass[border=5pt]{standalone}
% gnc-style.tex — 이 볼트의 TikZ 그림이 공유하는 스타일.
%
% 저널 그림 규약을 스타일 하나로 굳혀 둔다. 그림마다 선 굵기나 화살촉을
% 다시 정하지 않는다 — 그것이 그림이 제각각으로 보이는 원인이다.
%
% 쓰는 법:  % gnc-style.tex — 이 볼트의 TikZ 그림이 공유하는 스타일.
%
% 저널 그림 규약을 스타일 하나로 굳혀 둔다. 그림마다 선 굵기나 화살촉을
% 다시 정하지 않는다 — 그것이 그림이 제각각으로 보이는 원인이다.
%
% 쓰는 법:  % gnc-style.tex — 이 볼트의 TikZ 그림이 공유하는 스타일.
%
% 저널 그림 규약을 스타일 하나로 굳혀 둔다. 그림마다 선 굵기나 화살촉을
% 다시 정하지 않는다 — 그것이 그림이 제각각으로 보이는 원인이다.
%
% 쓰는 법:  \input{gnc-style}  (figures/src/ 안에서)

\usepackage{tikz}
\usepackage{amsmath}
\usepackage{xcolor}
\usetikzlibrary{arrows.meta, positioning, calc, fit, backgrounds, decorations.pathmorphing}

% ---- 색 --------------------------------------------------------------------
% 색은 강조 하나에만 쓴다. 나머지는 검정.
\definecolor{ink}{HTML}{1A1A1A}
\definecolor{accent}{HTML}{7C3AED}
\definecolor{warn}{HTML}{B91C1C}
\definecolor{soft}{HTML}{666666}

% ---- 선과 화살촉 -----------------------------------------------------------
\tikzset{
  % 신호선. 화살촉은 작고 뾰족한 것 하나뿐이다.
  sig/.style   = {ink, line width=0.5pt, -{Stealth[length=4pt, width=3pt]}},
  wire/.style  = {ink, line width=0.5pt},
  % 강조 경로 — 파선으로 구분한다. 흑백 인쇄에서도 살아야 하기 때문이다.
  asig/.style  = {accent, line width=0.5pt, densely dashed,
                  -{Stealth[length=4pt, width=3pt]}},
  awire/.style = {accent, line width=0.5pt, densely dashed},
  % 블록. 흰 바탕, 직각, 검은 테두리.
  blk/.style   = {draw=ink, line width=0.55pt, fill=white,
                  minimum width=17mm, minimum height=8mm, inner sep=2pt, font=\small},
  % 합산점
  sum/.style   = {draw=ink, line width=0.55pt, fill=white, circle,
                  inner sep=0pt, minimum size=4.6mm, font=\scriptsize},
  asum/.style  = {draw=accent, line width=0.55pt, fill=white, circle,
                  inner sep=0pt, minimum size=4.6mm, font=\scriptsize},
  % 분기점
  dot/.style   = {ink, fill=ink, circle, inner sep=0pt, minimum size=1.5mm},
  % 신호 이름 — 선 위에 작게
  sname/.style = {font=\small, inner sep=1.5pt},
  % 그림 안의 짧은 설명. 그림 안의 글은 최소로 둔다
  note/.style  = {font=\scriptsize, soft, inner sep=1.5pt},
  anote/.style = {font=\scriptsize, accent, inner sep=1.5pt},
  % 축
  axis/.style  = {ink, line width=0.5pt},
  tick/.style  = {font=\scriptsize, soft},
}

% 캡션 한 줄 — 그림 아래 가는 선과 함께
\newcommand{\gnccaption}[2]{%
  \node[anchor=north west, text width=#1, align=left, font=\scriptsize, ink]
        at (current bounding box.south west) {\vspace{2pt}#2};}
  (figures/src/ 안에서)

\usepackage{tikz}
\usepackage{amsmath}
\usepackage{xcolor}
\usetikzlibrary{arrows.meta, positioning, calc, fit, backgrounds, decorations.pathmorphing}

% ---- 색 --------------------------------------------------------------------
% 색은 강조 하나에만 쓴다. 나머지는 검정.
\definecolor{ink}{HTML}{1A1A1A}
\definecolor{accent}{HTML}{7C3AED}
\definecolor{warn}{HTML}{B91C1C}
\definecolor{soft}{HTML}{666666}

% ---- 선과 화살촉 -----------------------------------------------------------
\tikzset{
  % 신호선. 화살촉은 작고 뾰족한 것 하나뿐이다.
  sig/.style   = {ink, line width=0.5pt, -{Stealth[length=4pt, width=3pt]}},
  wire/.style  = {ink, line width=0.5pt},
  % 강조 경로 — 파선으로 구분한다. 흑백 인쇄에서도 살아야 하기 때문이다.
  asig/.style  = {accent, line width=0.5pt, densely dashed,
                  -{Stealth[length=4pt, width=3pt]}},
  awire/.style = {accent, line width=0.5pt, densely dashed},
  % 블록. 흰 바탕, 직각, 검은 테두리.
  blk/.style   = {draw=ink, line width=0.55pt, fill=white,
                  minimum width=17mm, minimum height=8mm, inner sep=2pt, font=\small},
  % 합산점
  sum/.style   = {draw=ink, line width=0.55pt, fill=white, circle,
                  inner sep=0pt, minimum size=4.6mm, font=\scriptsize},
  asum/.style  = {draw=accent, line width=0.55pt, fill=white, circle,
                  inner sep=0pt, minimum size=4.6mm, font=\scriptsize},
  % 분기점
  dot/.style   = {ink, fill=ink, circle, inner sep=0pt, minimum size=1.5mm},
  % 신호 이름 — 선 위에 작게
  sname/.style = {font=\small, inner sep=1.5pt},
  % 그림 안의 짧은 설명. 그림 안의 글은 최소로 둔다
  note/.style  = {font=\scriptsize, soft, inner sep=1.5pt},
  anote/.style = {font=\scriptsize, accent, inner sep=1.5pt},
  % 축
  axis/.style  = {ink, line width=0.5pt},
  tick/.style  = {font=\scriptsize, soft},
}

% 캡션 한 줄 — 그림 아래 가는 선과 함께
\newcommand{\gnccaption}[2]{%
  \node[anchor=north west, text width=#1, align=left, font=\scriptsize, ink]
        at (current bounding box.south west) {\vspace{2pt}#2};}
  (figures/src/ 안에서)

\usepackage{tikz}
\usepackage{amsmath}
\usepackage{xcolor}
\usetikzlibrary{arrows.meta, positioning, calc, fit, backgrounds, decorations.pathmorphing}

% ---- 색 --------------------------------------------------------------------
% 색은 강조 하나에만 쓴다. 나머지는 검정.
\definecolor{ink}{HTML}{1A1A1A}
\definecolor{accent}{HTML}{7C3AED}
\definecolor{warn}{HTML}{B91C1C}
\definecolor{soft}{HTML}{666666}

% ---- 선과 화살촉 -----------------------------------------------------------
\tikzset{
  % 신호선. 화살촉은 작고 뾰족한 것 하나뿐이다.
  sig/.style   = {ink, line width=0.5pt, -{Stealth[length=4pt, width=3pt]}},
  wire/.style  = {ink, line width=0.5pt},
  % 강조 경로 — 파선으로 구분한다. 흑백 인쇄에서도 살아야 하기 때문이다.
  asig/.style  = {accent, line width=0.5pt, densely dashed,
                  -{Stealth[length=4pt, width=3pt]}},
  awire/.style = {accent, line width=0.5pt, densely dashed},
  % 블록. 흰 바탕, 직각, 검은 테두리.
  blk/.style   = {draw=ink, line width=0.55pt, fill=white,
                  minimum width=17mm, minimum height=8mm, inner sep=2pt, font=\small},
  % 합산점
  sum/.style   = {draw=ink, line width=0.55pt, fill=white, circle,
                  inner sep=0pt, minimum size=4.6mm, font=\scriptsize},
  asum/.style  = {draw=accent, line width=0.55pt, fill=white, circle,
                  inner sep=0pt, minimum size=4.6mm, font=\scriptsize},
  % 분기점
  dot/.style   = {ink, fill=ink, circle, inner sep=0pt, minimum size=1.5mm},
  % 신호 이름 — 선 위에 작게
  sname/.style = {font=\small, inner sep=1.5pt},
  % 그림 안의 짧은 설명. 그림 안의 글은 최소로 둔다
  note/.style  = {font=\scriptsize, soft, inner sep=1.5pt},
  anote/.style = {font=\scriptsize, accent, inner sep=1.5pt},
  % 축
  axis/.style  = {ink, line width=0.5pt},
  tick/.style  = {font=\scriptsize, soft},
}

% 캡션 한 줄 — 그림 아래 가는 선과 함께
\newcommand{\gnccaption}[2]{%
  \node[anchor=north west, text width=#1, align=left, font=\scriptsize, ink]
        at (current bounding box.south west) {\vspace{2pt}#2};}


\begin{document}
\begin{tikzpicture}[x=1mm, y=1mm]

\def\Kd{2}
\def\XW{80}          % 가로 4 decade (0.1 .. 1000)
\def\YH{50}          % 세로 4.5 decade (0.1 .. ~3160)
\def\XA{20}          % 1 decade 당 mm  (=XW/4)
\def\YA{11.111}      % 1 decade 당 mm  (=YH/4.5)

% lx = log10(w), ly = log10(mag)  ->  화면 mm
\newcommand{\PX}[1]{{\XA*((#1)+1)}}
\newcommand{\PY}[1]{{\YA*((#1)+1)}}

% ---- 격자 -----------------------------------------------------------------
\foreach \e in {-1,0,1,2,3}
  \draw[soft, line width=0.2pt] (\PX{\e},0) -- (\PX{\e},\YH);
\foreach \e in {-1,0,1,2,3}
  \draw[soft, line width=0.2pt] (0,\PY{\e}) -- (\XW,\PY{\e});

% ---- 축 -------------------------------------------------------------------
\draw[axis] (0,0) -- (\XW,0) -- (\XW,\YH) -- (0,\YH) -- cycle;
\foreach \e/\lab in {-1/0.1, 0/1, 1/10, 2/100, 3/1000} {
  \draw[axis] (\PX{\e},0) -- (\PX{\e},-1.2);
  \node[tick, anchor=north] at (\PX{\e},-1.4) {\lab};
}
\foreach \e/\lab in {-1/0.1, 0/1, 1/10, 2/100, 3/1000} {
  \draw[axis] (0,\PY{\e}) -- (-1.2,\PY{\e});
  \node[tick, anchor=east] at (-1.6,\PY{\e}) {\lab};
}
\node[font=\small, anchor=north] at ({\XW/2},-6.5) {frequency $\omega$\ \ [rad/s]};
\node[font=\small, anchor=south, rotate=90] at (-11,{\YH/2})
     {gain applied at that frequency};

% ---- 이상 미분:  log10(Kd) + lx --------------------------------------------
\draw[ink, line width=0.9pt, domain=-1:3, samples=60, smooth]
      plot (\PX{\x}, \PY{log10(\Kd) + \x});

% ---- pseudo — 로그 공간에서 계산한다 ---------------------------------------
% |Kd N jw/(jw+N)| 의 log10 은
%     log10(Kd) + log10(N) + x - 0.5*log10(10^(2x) + N^2)
% 인데, pgfmath 는 고정소수라 10^6 에서 넘친다("Dimension too large").
% log10(10^a + 10^b) = max(a,b) + log10(1 + 10^(-|a-b|)) 로 바꾸면
% 중간값이 절대 커지지 않는다.  a = 2x, b = 2 log10(N).
% 바깥 중괄호를 붙이지 않는다 — \PY 가 이미 감싸므로 이중 중괄호가 되면
% pgfmath 가 "Missing number" 로 죽는다.
\newcommand{\LPS}[2]{log10(\Kd) + (#1) + (#2)
   - 0.5*( max(2*(#2), 2*(#1)) + log10(1 + pow(10, -abs(2*(#2) - 2*(#1)))) )}

\draw[ink, line width=0.7pt, densely dashed, domain=-1:3, samples=140, smooth]
      plot (\PX{\x}, \PY{\LPS{2}{\x}});

\draw[accent, line width=0.9pt, domain=-1:3, samples=140, smooth]
      plot (\PX{\x}, \PY{\LPS{1}{\x}});

% ---- 평탄부 표시 -----------------------------------------------------------
\draw[soft, line width=0.3pt, densely dotted]
      (\PX{1.4}, \PY{log10(200)}) -- (\XW, \PY{log10(200)});
\draw[soft, line width=0.3pt, densely dotted]
      (\PX{0.6}, \PY{log10(20)})  -- (\XW, \PY{log10(20)});

% 라벨은 곡선 위에 얹지 않는다. 이상미분은 끝점 왼쪽 위, 나머지 둘은
% 각자의 평탄선 위에 오른쪽 정렬로 붙인다.
\node[font=\scriptsize, ink, anchor=south east] at (\PX{2.62}, \PY{log10(\Kd*420)+0.10})
     {ideal $K_d s$};
\node[font=\scriptsize, ink, anchor=south east] at ({\XW-1}, \PY{log10(200)+0.13})
     {$N=100$};
\node[font=\scriptsize, accent, anchor=south east] at ({\XW-1}, \PY{log10(20)+0.13})
     {$N=10$};

% ---- 두 주파수 대역 --------------------------------------------------------
\draw[ink, line width=0.5pt] (\PX{log10(2)}, \YH) -- (\PX{log10(2)}, \YH+2.5);
\node[font=\scriptsize, anchor=south, align=center] at (\PX{log10(2)}, \YH+2.8)
     {the loop works here\\[-1pt]$\omega_n=2$};
\draw[ink, line width=0.5pt] (\PX{log10(314)}, \YH) -- (\PX{log10(314)}, \YH+2.5);
\node[font=\scriptsize, anchor=south, align=center] at (\PX{log10(314)}, \YH+2.8)
     {the noise reaches here\\[-1pt]$\pi/T_s=314$};

% ---- 측정값 ----------------------------------------------------------------
\node[font=\scriptsize, ink, anchor=north west, align=left, text width=44mm]
     at ({\XW+7}, \YH) {%
  The dotted lines are the plateaus,\\
  at $K_dN$: $200$ and $20$.\\[5pt]
  \textbf{At $\omega_n=2$} the three agree:\\
  $4.00$ against $3.92$.\\[3pt]
  \textbf{At $\pi/T_s=314$} they do not:\\
  $628$ against $20.0$ — a factor of $31.4$.\\[5pt]
  \textbf{Measured in section I}\\
  quiet-state actuator RMS\\
  ideal $8.75$, $N=10$ $0.23$\\
  peak demand $105.9$ against $4.26$};

\end{tikzpicture}
\end{document}
